\section{Evaluation framework}\label{sec:setup}
\subsection{Scenarios}
Missions are created by a seeded generator with the parameters in Table~\ref{tab:params}. The map is a square with a
side of 10\,km and the depot at its centre, and half of the tasks are clustered around three hotspots. Half of the tasks
are known at the start, and the others arrive in batches. Three fleet scales with 5, 10 and 20 UAVs and 25, 50 and 140
tasks are combined with three event profiles. The arrivals profile contains six arrival batches, the failures profile
three batches and two UAV failures, and the deadlines profile three batches and four tightened deadlines. A fourth scale
with 40 UAVs and 280 tasks is used for scalability only.

\begin{table}[t]
\caption{Scenario generator parameters (distances in m, times in s, energy in metres of flight).}\label{tab:params}
\centering\small
\setlength{\tabcolsep}{4pt}
\begin{tabular}{ll@{\qquad}ll}
\toprule
Map side / horizon & 10\,000 / 1\,800 & Service time & $U[30,90]$ \\
UAV speed & $U[10,16]$\,m/s & Time-window slack & $U[180,600]$ \\
Energy capacity & $U[18,28]$\,km & Release delay & $U[0,120]$ \\
Energy reserve & 10\% of capacity & Rewards 1, 2, 3, 5 & probabilities .4, .3, .2, .1 \\
Hover cost & 6 per second & Hotspots & 3, std.\ 600, 50\% of tasks \\
Tasks known at start & 50\% & Scales (UAVs/tasks) & 5/25, 10/50, 20/140, 40/280 \\
\bottomrule
\end{tabular}
\end{table}

\subsection{Static snapshots with exact execution}
The stopping rules are first compared without execution uncertainty. Each mission is simulated once with a
deterministic reference planner, and every replanning call after the start is stored with its previous plan. All
methods are evaluated on these snapshots, and regret is measured against the best known objective of each snapshot,
which includes a 3\,s ALNS run and, for the small and medium scales, the OR-Tools routing solver~\cite{ortools} with
guided local search~\cite{voudouris1999gls}. Because all ALNS variants follow the same search trajectory, whole families
of each rule are evaluated exactly offline, with 12 fixed budgets, 10 stagnation patience values and 9 compute prices,
and each baseline is compared with AnytimeSwarm at the same mean computation. This benchmark uses its own pre-registered
test seeds 2000--2029 (\val{NsnapTest} snapshots of 270 missions), and it was reproduced on the final software platform
before the execution-aware study (Section~\ref{sec:repro}).

\subsection{Tier B: kinematic execution with disturbances and latency}
Tier B executes each published plan in continuous mission time. The duration of a leg follows a trapezoidal speed
profile with a horizontal acceleration limit. The cruise airspeed of each (UAV, task) leg deviates from the nominal
speed by a multiplicative error, and a constant wind of fixed magnitude and random direction changes the ground speed
through the wind triangle. Service durations are multiplied by a random factor. Energy is counted as metres flown
through the air plus the hover cost. The planner never sees these disturbances. It receives the state of each UAV as
telemetry, and the remaining time of a committed task is estimated with the nominal model from the actual position of
the UAV. The dissemination latency of a plan is lognormal, and each message is lost with a given probability and resent
after a timeout. Table~\ref{tab:levels} lists the disturbance and latency levels. They span a range from mild to severe
conditions and are not claimed to represent a particular vehicle.

All random draws use common random numbers~\cite{law2015simulation}. The speed error of a leg depends only on the
mission seed, the UAV and the task, the service factor only on the mission seed and the task, the wind only on the
mission seed, and the latency only on the mission seed and the event. Two planners that assign the same task to the same
UAV therefore experience the same disturbance, independently of the number of replanning calls. With zero disturbance,
zero latency and the event trigger, tier B reproduces the discrete-event simulator of the original formulation exactly,
which is verified by unit tests.

\begin{table}[t]
\caption{Disturbance and latency levels of tier B.}\label{tab:levels}
\centering\small
\setlength{\tabcolsep}{4pt}
\begin{tabular}{lcccc@{\quad}lccc}
\toprule
Execution & accel. (m/s$^2$) & speed c.v. & wind (m/s) & service c.v. & Latency & median (s) & $\sigma_{\log}$ & loss / timeout \\
\midrule
nominal & $\infty$ & 0 & 0 & 0 & none & 0 & -- & -- \\
low & 3.0 & 0.03 & 2 & 0.05 & low & 0.2 & 0.5 & -- \\
mid & 2.0 & 0.07 & 5 & 0.15 & high & 2.0 & 0.8 & 5\% / 5\,s \\
high & 1.5 & 0.12 & 8 & 0.25 & & & & \\
\bottomrule
\end{tabular}
\end{table}

\subsection{Tier C: PX4 software-in-the-loop with Gazebo and ROS 2}
Tier C flies missions with three X500 quadrotors, each controlled by its own PX4 v1.17 autopilot in software-in-the-loop
mode~\cite{meier2015px4}, in Gazebo Harmonic~\cite{koenig2004gazebo}, on ROS~2 Jazzy~\cite{macenski2022ros2}
(Fig.~\ref{fig:tierc}). A uXRCE-DDS bridge exposes the internal topics of each autopilot to ROS~2. Three nodes form the
mission stack. The mission manager runs the scenario timeline, builds the planning snapshot from the fleet state, calls
the planner and releases the plan. The executor streams offboard setpoints to every autopilot at 10\,Hz, so that PX4
does not leave offboard mode while the planner computes. It flies each leg with a carrot that moves along the straight
leg at the commanded ground speed, holds the UAV at the task for the service duration and reports service start, finish
and skips. A visualisation node draws tasks, routes and planning decisions in RViz~2 and as coloured pillars in the
Gazebo scene. Planning time is measured with a monotonic wall clock, and every mission time is simulation time.

A tier C mission is a tier B scenario with three UAVs and fifteen tasks, which keeps the ratio of five tasks per UAV of
the small scale, rescaled by $S_X=0.1$ in length and $S_T=0.25$ in time. Speeds are multiplied by $S_X/S_T$, energy by
$S_X$, the hover cost by $S_X/S_T$ and the distance price by $1/S_X$. Every dimensionless ratio of the decision problem
is preserved, so a tier C mission is the same decision problem as its tier B counterpart, flown in a 1\,km world in
450\,s with cruise speeds of 4--6.4\,m/s. A planning delay of $\tau$ in tier C corresponds to $\tau/S_T$ in tier B, so
the latency factors are related by $\kappa_C=S_T\kappa_B$. Disturbances use the same common random numbers as tier B.
The commanded ground speed of each leg follows the speed error and the wind triangle, with the wind scaled by
$S_X/S_T$, and the service duration follows the same factor. Acceleration, attitude dynamics, position tracking and
communication are those of the simulated vehicles. Gazebo runs in lockstep with PX4 at twice real time, which
requires $\kappa_C\ge2$ so that the emulated planning delay dominates the real one.

\begin{figure}[t]
\centering
\fig{fig_tierc}{0.95\linewidth}
\caption{Tier C architecture. The planner and the mission timeline run in the mission manager; the executor streams
offboard setpoints to three PX4 autopilots that fly X500 quadrotors in Gazebo.}\label{fig:tierc}
\end{figure}

\subsection{Compared methods}
The constructive baseline is regret-2 repair, that is, FFR without search. ALNS is run with the full budget of
250\,ms, with a fixed budget of 40\,ms, which equals the mean computation of AnytimeSwarm on development data, and with
a stagnation rule that stops after 20\% of the budget without improvement. AnytimeSwarm is used exactly as frozen after
development, without retraining, and an ablation removes its delay cost. AnytimeSwarm-R uses the error model of the
execution level of the mission and the expected latency of its latency level. A further ablation sets its error model to
zero, so that only the latency term remains, and a sensitivity study uses the error model of another level. For 40
UAVs, the fixed budget matched to the computation of AnytimeSwarm on development data is 180\,ms. All ALNS variants
share operators, parameters, initial plan and seed, and every call has a hard budget of 250\,ms.

\subsection{Protocol and statistics}\label{sec:protocol}
Development used seeds 0--9 in tier B and 0--2 in tier C. The predictor was trained on development trajectories, and
the compute price $\mu=10^{-5}$ was selected by a rule stated before any test run. The tier B error models were fitted
on 90 development missions per level flown by ALNS with 40\,ms, and the tier C models on nine development flights per
level. The fitted models are listed in Table~\ref{tab:errmodels}. The experiment configurations, the error models, the
predictor and the experiment code were then frozen with SHA-256 hashes, and the held-out runs checked these hashes
before starting. Tier B uses held-out seeds 3000--3029, which gives 270 missions for the three scales and three
profiles, and 90 missions for 40 UAVs. Tier C uses held-out seeds 4000--4009 with three profiles, that is, 30 missions
per method. The primary tier B matrix crosses two execution levels (low, mid), two latency levels (none, high), two
latency factors ($\kappa_B=10, 50$) and two triggers (event, periodic every 60\,s).

The primary hypothesis is that AnytimeSwarm-R serves more reward than ALNS with 40\,ms, AnytimeSwarm and the stagnation
rule over the primary matrix. For each mission, the utility difference is averaged over the sixteen configurations of
the matrix, and the mean over missions is tested with a paired bootstrap~\cite{efron1993bootstrap} of 10\,000 resamples.
The three comparisons are corrected with the Holm procedure~\cite{holm1979simple}. The same pooled test compares the
computation per planning call with that of AnytimeSwarm and the periodic trigger with the event trigger. In tier C,
AnytimeSwarm-R is compared with ALNS with 40\,ms and with AnytimeSwarm, again with Holm correction, and the same 30
missions are also run in tier B to assess whether tier B ranks the methods as tier C does. All other comparisons are
descriptive. A tier C run in which a vehicle does not reach its altitude within 90\,s of simulation time is aborted and
repeated once with the same seed. The full protocol was written before the held-out runs and is available with the code,
together with a list of the deviations from it. The only deviation that affects data concerns four tier C runs with a
low real-time factor (Section~\ref{sec:res-tierc}).
Tier B runs four missions in parallel on an Intel Core i7-11800H laptop, with all methods of a configuration run in
turn, so that they share the same machine load. All runs use Python 3.12 and pinned package versions. Every plan is
checked for feasibility before publication, and a separately written validator that shares no code with the planner
found no infeasible plan among \val{JValidatorPlans} development plans of all methods.

\inputtable{tab_errmodels}
